\Article{5}
Der erste Beweis des Satzes stammt von dem ausgezeichneten Mathematiker \textsc{d'Alembert}, \textit{Recherches sur le calcul intégral, Histoire de l'Acad. de Berlin, Année} 1746, S. 182 ff. Er findet sich auch bei \textsc{Bougainville}, \textit{Traité du calcul intégral, à Paris} 1754, S. 47 ff. Die Hauptschritte dieser Methode sind folgende.

Zunächst zeigt er: Wird irgendeine Funktion $X$ der veränderlichen Größe $x$ entweder für $x=0$ oder für $x=\infty$ gleich $=0$ und kann sie einen unendlich kleinen positiven reellen Wert annehmen, wenn man $x$ einen reellen Wert gibt, so kann diese Funktion auch einen unendlich kleinen negativen reellen Wert durch einen Wert von $x$ annehmen, der entweder reell ist oder die imaginäre Form $p+q\sqrt{-1}$ hat. Bezeichnet nämlich $\Omega$ einen unendlich kleinen Wert von $X$ und $\omega$ den entsprechenden Wert von $x$, so behauptet er, $\omega$ lasse sich durch eine rasch konvergierende Reihe $a\Omega^\alpha+b\Omega^\beta+c\Omega^\gamma$ usw. ausdrücken, deren Exponenten $\alpha$, $\beta$, $\gamma$ usw. fortwährend zunehmende rationale Größen sind, die daher wenigstens von einer bestimmten Entfernung vom Anfang an positiv werden und die Glieder, in denen sie vorkommen, unendlich klein machen. Kommt nun unter diesen Exponenten kein Bruch mit geradem Nenner vor, so werden alle Glieder der Reihe sowohl für einen positiven als auch für einen negativen Wert von $\Omega$ reell. Finden sich dagegen unter jenen Exponenten Brüche mit geradem Nenner, so haben für einen negativen Wert von $\Omega$ die entsprechenden Glieder bekanntlich die Form $p+q\sqrt{-1}$. Wegen der unendlichen Konvergenz der Reihe genügt es aber im ersten Fall, nur das erste (das heißt größte) Glied beizubehalten; im zweiten braucht man nicht über dasjenige Glied hinauszugehen, das zuerst einen imaginären Teil hervorbringt.

Durch ähnliche Überlegungen lasse sich zeigen: Kann $X$ durch einen reellen Wert von $x$ einen unendlich kleinen negativen Wert annehmen, so kann diese Funktion durch einen reellen Wert von $x$ oder einen imaginären der Form $p+q\sqrt{-1}$ einen unendlich kleinen positiven reellen Wert erhalten.

Daraus schließt er zweitens, dass es auch einen endlichen reellen Wert von $X$ gibt, im ersten Fall einen negativen, im zweiten einen positiven, der durch einen imaginären Wert von $x$ der Form $p+q\sqrt{-1}$ hervorgebracht werden kann.

\SourceLeaf{013}
Hieraus folgt: Ist $X$ eine solche Funktion von $x$, die durch einen reellen Wert $v$ von $x$ einen reellen Wert $V$ erhält und durch einen reellen Wert von $x$ auch einen um eine unendlich kleine Größe größeren oder kleineren reellen Wert annimmt, so kann sie auch einen reellen Wert erhalten, der entsprechend um eine unendlich kleine und folglich um eine endliche Größe kleiner oder größer als $V$ ist, indem man $x$ einen Wert der Form $p+q\sqrt{-1}$ gibt. Das ergibt sich ohne Schwierigkeit aus dem Vorhergehenden, wenn man sich für $X$ den Ausdruck $V+Y$ und für $x$ den Ausdruck $v+y$ eingesetzt denkt.

Schließlich behauptet der ausgezeichnete \textsc{d'Alembert}: Nimmt man an, dass $X$ ein ganzes Intervall zwischen zwei reellen Werten $R,S$ durchlaufen kann (das heißt sowohl $R$ als auch $S$ und allen dazwischenliegenden reellen Werten gleich werden kann), indem man $x$ stets Werte der Form $p+q\sqrt{-1}$ gibt, so kann die Funktion $X$ noch um jede endliche reelle Größe vermehrt oder vermindert werden (je nachdem, ob $S>R$ oder $S<R$ ist), wobei $x$ stets die Form $p+q\sqrt{-1}$ behält. Gäbe es nämlich eine reelle Größe $U$ (zwischen der und $R$ der Wert $S$ liegen soll), der $X$ durch einen solchen Wert von $x$ nicht gleich werden könnte, so müsste es einen \textit{größten} Wert von $X$ geben (wenn $S>R$ ist; einen kleinsten dagegen, wenn $S<R$ ist), etwa $T$, den sie durch einen Wert $p+q\sqrt{-1}$ von $x$ erreichte, sodass man $x$ keinen Wert derselben Form geben könnte, der die Funktion $X$ auch nur um den kleinsten Betrag näher an $U$ brächte. Setzt man nun in dieser Gleichung zwischen $X$ und $x$ überall $p+q\sqrt{-1}$ für $x$ ein und setzt dann sowohl den reellen Teil als auch den Teil, der den Faktor $\sqrt{-1}$ enthält, nach Weglassung dieses Faktors gleich null, so kann man aus den beiden entstehenden Gleichungen (in denen $p,q$ und $X$ zusammen mit Konstanten vorkommen) durch Elimination zwei andere gewinnen. In der einen stehen $p,X$ und Konstanten; die andere ist von $p$ frei und enthält nur $q,X$ und Konstanten. Da also $X$ durch reelle Werte von $p,q$ alle Werte von $R$ bis $T$ durchlaufen hat, kann $X$ nach dem Vorhergehenden noch näher an den Wert $U$ herankommen, wenn man $p,q$ entsprechend Werte der Formen $\alpha+\gamma\sqrt{-1}$ und $\beta+\delta\sqrt{-1}$ gibt. Dann wird aber $x=\alpha-\delta+(\gamma+\beta)\sqrt{-1}$, hat also weiterhin die Form $p+q\sqrt{-1}$, entgegen der Voraussetzung.

Nimmt man jetzt an, dass $X$ eine Funktion wie $x^m+Ax^{m-1}+Bx^{m-2}+\text{usw.}+M$ bezeichnet, so sieht man ohne Schwierigkeit, dass man $x$ solche reellen Werte geben kann, dass $X$ ein ganzes Intervall zwischen zwei reellen Werten durchläuft. Daher kann $x$ einen Wert der Form $p+q\sqrt{-1}$ erhalten, durch den $X$ gleich $=0$ wird. Q. E. D.\HistoricalNote{2}
\SourceLeaf{014}
\Article{6}
Die Einwände, die sich gegen den Beweis von \textsc{d'Alembert} erheben lassen, laufen im Wesentlichen auf Folgendes hinaus.

1. Der ausgezeichnete d'A. zweifelt nicht an der \textit{Existenz} der Werte von $x$, denen gegebene Werte von $X$ entsprechen, sondern setzt sie voraus und untersucht nur die \textit{Form} jener Werte.

Obwohl dieser Einwand für sich genommen sehr schwer wiegt, betrifft er hier doch nur die Ausdrucksweise, die sich leicht so ändern lässt, dass er vollständig entfällt.

2. Die Behauptung, $\omega$ lasse sich stets durch eine Reihe der von ihm angegebenen Art ausdrücken, ist gewiss falsch, wenn $X$ auch eine beliebige transzendente Funktion bezeichnen soll (wie d'A. an mehreren Stellen andeutet). Das ist zum Beispiel klar, wenn man $X=e^{\frac1x}$ oder $x=\frac1{\log X}$ setzt. Beschränken wir jedoch den Beweis auf den Fall, dass $X$ eine algebraische Funktion von $x$ ist (was für die vorliegende Aufgabe genügt), so ist der Satz allerdings richtig. --- Im Übrigen hat d'A. zur Bestätigung seiner Voraussetzung nichts angeführt. Der berühmte \textsc{Bougainville} setzt voraus, dass $X$ eine algebraische Funktion von $x$ ist, und empfiehlt zur Auffindung der Reihe das \textsc{Newton}sche Parallelogramm.

3. Er bedient sich unendlich kleiner Größen freier, als es mit geometrischer Strenge vereinbar ist oder wenigstens in unserer Zeit (in der sie mit Recht in schlechtem Ruf stehen) von einem gewissenhaften Analysten zugestanden würde. Auch erläutert er den Sprung von einem unendlich kleinen Wert von $\Omega$ zu einem endlichen nicht deutlich genug. Seinen Satz, dass $\Omega$ auch einen endlichen Wert annehmen kann, scheint er weniger aus der Möglichkeit eines unendlich kleinen Wertes von $\Omega$ zu folgern als vielmehr daraus: Bezeichnet $\Omega$ eine sehr kleine Größe, so komme man wegen der raschen Konvergenz der Reihe dem wahren Wert von $\omega$ umso näher, je mehr Glieder der Reihe man nimmt; oder je mehr Glieder man in die für $\omega$ genommene Summe einbezieht, desto genauer werde die Gleichung erfüllt, die die Beziehung zwischen $\omega$ und $\Omega$ beziehungsweise $x$ und $X$ ausdrückt. Abgesehen davon, dass diese ganze Argumentation zu unbestimmt erscheint, um daraus irgendeinen strengen Schluss zu ziehen, bemerke ich, dass es durchaus Reihen gibt, die, wie klein man auch den Wert der Größe ansetzt,
\SourceLeaf{015}
nach deren Potenzen sie fortschreiten, dennoch immer divergieren, sodass man, wenn man sie nur weit genug fortsetzt, zu Gliedern gelangen kann, die größer als jede gegebene Größe sind\HistoricalNote{3}. Das tritt ein, wenn die Koeffizienten der Reihe eine hypergeometrische Progression bilden. Daher hätte notwendig bewiesen werden müssen, dass im vorliegenden Fall eine solche hypergeometrische Reihe nicht entstehen kann.

Im Übrigen scheint mir, dass der ausgezeichnete d'A. hier zu Unrecht zu unendlichen Reihen gegriffen hat und dass diese zur Begründung dieses Fundamentalsatzes der Gleichungslehre nicht geeignet sind.

4. Aus der Annahme, dass $X$ den Wert $S$, nicht aber den Wert $U$ annehmen kann, folgt noch nicht, dass zwischen $S$ und $U$ notwendig ein Wert $T$ liegt, den $X$ erreichen, aber nicht überschreiten kann. Es bleibt noch ein anderer Fall: Zwischen $S$ und $U$ könnte eine Grenze liegen, der $X$ beliebig nahe kommen kann, ohne sie jemals zu erreichen. Aus den vom ausgezeichneten d'A. angeführten Gründen folgt nur, dass $X$ jeden einmal erreichten Wert noch um eine endliche Größe überschreiten kann. Ist sie etwa gleich $=S$ geworden, so kann sie noch um eine endliche Größe $\Omega$ vermehrt werden; danach kommt ein neuer Zuwachs $\Omega'$ hinzu, dann wieder ein Zuwachs $\Omega''$ usw., sodass, wie viele Zuwächse auch schon hinzugefügt sind, keiner als der letzte angesehen werden darf, sondern stets noch ein neuer hinzukommen kann. Doch obwohl die \textit{Anzahl} der möglichen Zuwächse durch keine Grenzen beschränkt ist, könnte es durchaus geschehen, dass bei fortwährender Abnahme der Zuwächse $\Omega,\Omega',\Omega''$ usw. die Summe $S+\Omega+\Omega'+\Omega''$ usw. dennoch eine Grenze niemals erreicht, wie viele Glieder man auch betrachtet.

Dieser Fall kann zwar nicht vorkommen, wenn $X$ ein Polynom in $x$ bezeichnet; ohne Beweis, dass dies nicht geschehen kann, muss die Methode dennoch notwendig als unvollständig gelten. Ist dagegen $X$ eine transzendente oder auch eine gebrochene algebraische Funktion, so kann jener Fall durchaus eintreten, zum Beispiel immer dann, wenn einem bestimmten Wert von $X$ ein unendlich großer Wert von $x$ entspricht\SourceLeaf{016}. Dann scheint sich die Methode von \textsc{d'Alembert} nicht ohne viele Umwege und in manchen Fällen vielleicht überhaupt nicht auf unbezweifelbare Grundlagen zurückführen zu lassen.

Aus diesen Gründen kann ich den Beweis von \textsc{d'Alembert} nicht für befriedigend halten. Dennoch scheint mir seine eigentliche Beweiskraft durch alle diese Einwände keineswegs gebrochen. Ich glaube, dass sich auf derselben Grundlage (wenn auch auf ganz andere Weise und wenigstens mit größerer Umsicht) nicht nur ein strenger Beweis unseres Satzes aufbauen lässt, sondern daraus auch alles gewonnen werden kann, was man sich für die Theorie \textit{transzendenter} Gleichungen wünschen könnte. Über diese sehr wichtige Sache werde ich bei anderer Gelegenheit ausführlicher sprechen; man vergleiche inzwischen unten Artikel 24.
